\documentclass[final]{agujournal2025}

\usepackage{subfiles}
\let\citep\cite

\usepackage{graphicx}
\usepackage{booktabs}
\usepackage{amsmath}


\begin{document}

\journalname{Earth's Future}

% Your title can be multiple lines long.
\title{Susceptibility to Abrupt Thaw across Alaska's Permafrost Landscapes}

\authors{
Ethan Pierce\affil{1,2,3},
Irina Overeem\affil{2,3},
Hailey Webb\affil{4,5},
Kevin Rozmiarek\affil{2,3},
Merritt Turetsky\affil{4,5}
}

\affiliation{1}{Thayer School of Engineering, Dartmouth College, Hanover, NH, USA}
\affiliation{2}{Institute of Arctic and Alpine Research, University of Colorado Boulder, Boulder, CO, USA}
\affiliation{3}{Department of Earth Science, University of Colorado Boulder, Boulder, CO, USA}
\affiliation{4}{Renewable and Sustainable Energy Institute, University of Colorado Boulder, Boulder, CO, USA}
\affiliation{5}{Ecology and Evolutionary Biology, University of Colorado Boulder, Boulder, CO, USA}

\correspondingauthor{Ethan Pierce}{ethan.g.pierce@dartmouth.edu}

\authorrunninghead{Ignored.}
\titlerunninghead{Ignored.}

\keypoints
% {Each must be 140 characters or fewer with no special characters or punctuation and must be complete sentences}
{We train a machine learning model to classify abrupt and non-abrupt thaw across Alaska's permafrost region}
{Abrupt thaw is favored over non-abrupt thaw across roughly one quarter of the applicable permafrost area}
{Shapley (SHAP) analysis identifies terrain characteristics and mean annual temperature regime as the leading indicators of thaw mode}

\maketitle

\subfile{sections/abstract}
\newpage

\subfile{sections/intro}

\subfile{sections/data}

\subfile{sections/methods}

\subfile{sections/results}

\subfile{sections/discussion}

\subfile{sections/conclusion}

\newpage
% Your Data Availability Statement is an unnumbered section right before the bibliography.
\section*{Open Research Statement}
\noindent The statewide susceptibility map, its area-of-applicability layer, the prediction datacube, and the training feature table from this study are archived on Zenodo \citep{pierce2026data} under a CC BY 4.0 license. The code that builds the feature table and prediction datacube, trains and evaluates the classifier, computes SHAP values, and produces the map and figures is archived on Zenodo \citep{pierce2026code} as release v1.0.0 of the project repository, together with the trained model, under the GNU GPL v3.0 license. The analysis relies on open-source, community software, including XGBoost \citep{chen2016xgboost}, TreeSHAP \citep{lundberg2020local}, scikit-learn \citep{pedregosa2011scikit}, and xarray \citep{hoyer2017xarray}. Thaw mode labels come from the Alaska Permafrost Thaw Database v2.0.0 \citep{webb2026thawdb_data}, which is released under CC0 and included in the code repository. All other inputs are publicly available from their providers and are not redistributed here. We accessed the USGS 3DEP 10 m DEM \citep{usgs3dep}, MERIT Hydro v1.0.1 \citep{merit_hydro}, MODIS MCD64A1 v6.1 \citep{modis_mcd64a1_061}, SoilGrids 2.0 \citep{poggio2021soilgrids}, WorldClim v1.4 \citep{worldclim_v14}, and Daymet V4 \citep{daymet_v4} through Google Earth Engine, and the code lists each asset identifier. We downloaded NLCD 2016 for Alaska \citep{dewitz2019national}, ALFRESCO model outputs \citep{alfresco}, IRYP v2 \citep{iryp_v2}, the permafrost probability map of \citeA{obu2018pangaea}, the Circumpolar Thermokarst Landscapes map \citep{olefeldt2015ornl}, and the Level III Ecoregions of Alaska \citep{epa_ecoregions_ak} directly from their providers.

\section*{Conflict of Interest declaration}
\noindent The authors declare there are no conflicts of interest for this manuscript.

\acknowledgments
The research reported here was funded in part by the Army Research Office/Army Research Laboratory via grant \#W911NF-23-1-0311 to the University of Colorado Boulder.  Any errors and opinions are not those of the Army Research Office and are attributable solely to the author(s). This material is also based upon work supported by the Broad Agency Announcement Program and the Cold Regions Research and Engineering Laboratory (ERDC-CRREL) under Contract No. W913E524C0003.

% Bibliography
%\cite{*} % If you have uncited bibliography items.
\bibliography{bib} % Replace ``wiley'' with your bibliography's file name.  Do not specify its extension, e.g., .bib.

\appendix  % Put this first.

% Can repeat for each appendix.
\newpage
\section{Data sources}\label{appx:data-sources}
\subfile{tables/data-sources}
\clearpage

\newpage
\section{SHAP feature families}\label{appx:feature-families}
\subfile{appendices/feature-families}
\clearpage

\end{document}
